\section*{Capítulo \ref{ch:b3:topology} --- Topología general}

\begin{solution}{exo:b3:topology:1}
(a) Cuatro \omterm{def:b3:topology:topology}{topologías} sobre $\{a,b\}$: la indiscreta
$\{\varnothing, X\}$; la discreta; y las dos \omterm{def:b3:topology:topology}{topologías} de
\emph{Sierpiński} $\{\varnothing, \{a\}, X\}$ y $\{\varnothing,
\{b\}, X\}$. Las dos últimas son
\omterm{def:b3:topology:continuity}{homeomorfas} (intercámbiese $a, b$):
tres clases. \omterm{def:b3:topology:connected}{Conexas}: todas salvo la
discreta (solo la \omterm{def:b3:topology:topology}{topología} discreta
contiene un conjunto propio no vacío \omterm{def:b3:topology:topology}{abierto} y cerrado a la vez).
\omterm{def:b3:topology:hausdorff}{De Hausdorff}: solo la discreta (en
las demás, el único \omterm{def:b3:topology:topology}{entorno} de $b$ es
$X$, o el de $a$).

(b) $\Q$: \omterm{def:b3:topology:interior}{interior} $\varnothing$
(todo intervalo contiene irracionales),
\omterm{def:b3:topology:interior}{clausura} $\R$ (densidad), frontera
$\R$. $[0,1)
\cup \{2\}$: \omterm{def:b3:topology:interior}{interior} $(0,1)$,
\omterm{def:b3:topology:interior}{clausura} $[0,1]\cup\{2\}$, frontera
$\{0, 1, 2\}$. $\{1/n\}$: \omterm{def:b3:topology:interior}{interior}
$\varnothing$, \omterm{def:b3:topology:interior}{clausura}
$\{1/n\}\cup\{0\}$, frontera la
\omterm{def:b3:topology:interior}{clausura} misma.
\end{solution}

\begin{solution}{exo:b3:topology:2}
(a) Todo \omterm{def:b3:topology:topology}{abierto} $V \subseteq Y$ es una unión $\bigcup B_j$ de conjuntos
de la \omterm{def:b3:topology:basis}{base}, y
$f^{-1}(V) = \bigcup f^{-1}(B_j)$: si estos últimos son \omterm{def:b3:topology:topology}{abiertos},
también lo es el primero; el recíproco es trivial.

(b) $f^{-1}\bigl(\intoo{\frac12}{\frac32}\bigr) = [0, \infty)$
no es \omterm{def:b3:topology:topology}{abierto}: $f$ no es \omterm{def:b3:topology:continuity}{continua};
la \omterm{def:b3:topology:continuity}{continuidad} por la derecha en
cada punto es clara ($f$ es localmente constante a la derecha de todo
punto). Los conjuntos $[a, b)$ cumplen el criterio de
\omterm{def:b3:topology:basis}{base} (la~\cref{def:b3:topology:basis}):
$[a,b)\cap[c,d) = [\max(a,c),
\min(b,d))$. La \omterm{def:b3:topology:topology}{topología} de
Sorgenfrey es más fina que la usual, pues
$(a, b) = \bigcup_{a < c < b}[c, b)$; y estrictamente: $[0,1)$ es
\omterm{def:b3:topology:topology}{abierto} de Sorgenfrey pero no \omterm{def:b3:topology:topology}{abierto} usual.
\omterm{def:b3:topology:continuity}{Continuidad} desde la recta de
Sorgenfrey en $x$: un \omterm{def:b3:topology:topology}{entorno} de
Sorgenfrey de $x$ contiene un conjunto de la
\omterm{def:b3:topology:basis}{base} $[a, b) \ni x$ y, por tanto,
contiene $[x, b)$ --- de modo que la condición de
\omterm{def:b3:topology:continuity}{continuidad} se lee: para todo
$\varepsilon > 0$ existe $\delta > 0$ con
$\abs{f(t) - f(x)} < \varepsilon$ siempre que $x \leq t < x + \delta$.
Eso es exactamente la \omterm{def:b3:topology:continuity}{continuidad}
por la derecha en $x$.
\end{solution}

\begin{solution}{exo:b3:topology:3}
Dos \omterm{def:b3:topology:topology}{abiertos} no vacíos tienen complementario finito, luego su
intersección también lo tiene: es no vacía (el conjunto es infinito) ---
no hay dos puntos con \omterm{def:b3:topology:topology}{entornos}
disjuntos: no es \omterm{def:b3:topology:hausdorff}{de Hausdorff} y, por
tanto, no es metrizable (la~\cref{def:b3:topology:hausdorff}). Sea
$(x_n)$ inyectiva y sea $x$ arbitrario: un
\omterm{def:b3:topology:topology}{entorno} de $x$ contiene un \omterm{def:b3:topology:topology}{abierto}
cofinito $U$; los finitos puntos de $X\setminus U$ son alcanzados por un
número finito de índices a lo sumo (inyectividad), de modo que una cola
de la sucesión está en $U$: $x_n \to x$, para \emph{todo} $x$. La
demostración de la unicidad necesita dos
\omterm{def:b3:topology:topology}{entornos} disjuntos para separar dos
supuestos límites --- precisamente lo que aquí falla.
\end{solution}

\begin{solution}{exo:b3:topology:4}
(a) La \omterm{def:b3:topology:continuity}{continuidad} es por
construcción (la~\cref{def:b3:topology:constructions}). Que sea abierta:
un \omterm{def:b3:topology:topology}{abierto} $W
\subseteq X \times Y$ es unión de cajas $U_i\times V_i$, y
$\pi_X(W) = \bigcup U_i$ (las cajas no vacías se proyectan sobre sus
factores), que es \omterm{def:b3:topology:topology}{abierto}. No cerrada: $H = \{(x,y) : xy = 1\}$ es
cerrado en $\R^2$ (preimagen de $\{1\}$ por el producto, que es
\omterm{def:b3:topology:continuity}{continuo}), pero
$\pi_X(H) = \R\setminus\{0\}$ no es cerrado.

(b) ($\Rightarrow$) Las proyecciones son
\omterm{def:b3:topology:continuity}{continuas}. ($\Leftarrow$) Sea
$x^{(k)} \to x$ componente a componente y sea $W = \prod U_n$ un
\omterm{def:b3:topology:topology}{entorno} de la
\omterm{def:b3:topology:basis}{base} de $x$, con $U_n = X_n$ salvo para
$n \in F$ finito. Para cada $n \in F$, $x_n^{(k)} \in U_n$ si $k \geq
K_n$; para $k \geq \max_F K_n$, $x^{(k)} \in W$. La fórmula $d$ define
una métrica (cada sumando lo es, salvo la verificación estándar de que
$\min(1, d_n)$ lo es); sus bolas: $B_d(x,
2^{-N})$ contiene la caja $\{y : d_n(x_n,y_n) < 2^{-N},\
n \leq N'\}$ de la \omterm{def:b3:topology:basis}{base} para $N'$
adecuados (la cola $\sum_{n > N'} 2^{-n}$ es pequeña) y, recíprocamente,
toda caja de la \omterm{def:b3:topology:basis}{base} contiene una
$d$-bola: las dos \omterm{def:b3:topology:topology}{topologías} tienen los mismos
\omterm{def:b3:topology:topology}{entornos} de cada punto.

(c) En la \omterm{def:b3:topology:constructions}{topología producto},
$x^{(k)} \to 0$ por (b). El \omterm{def:b3:topology:topology}{abierto} de cajas
$U = \prod_n \intoo{-\tfrac1n}{\tfrac1n}$ contiene $0$, pero
$x^{(k)} \notin U$ para todo $k$ (la coordenada $n$-ésima vale $1/k
\geq 1/n$ cuando $n \geq k$): ninguna cola entra en $U$. La
\omterm{def:b3:topology:topology}{topología} de cajas es estrictamente
más fina y no es una \omterm{def:b3:topology:topology}{topología} de
tipo producto a efectos de convergencia.
\end{solution}

\begin{solution}{exo:b3:topology:5}
$\varphi \colon \R \to S^1$, $x \mapsto \eu^{2\iu\pi x}$, es
\omterm{def:b3:topology:continuity}{continua}, sobreyectiva y constante
en las clases módulo $\Z$: induce una biyección
\omterm{def:b3:topology:continuity}{continua}
$\bar\varphi \colon \R/\Z \to
S^1$ (la~\cref{prop:b3:topology:universal}(b)). La proyección $\pi$ es
\emph{abierta}: para $U \subseteq \R$ \omterm{def:b3:topology:topology}{abierto},
$\pi^{-1}(\pi(U)) = \bigcup_{n}(U + n)$ es \omterm{def:b3:topology:topology}{abierto}, luego $\pi(U)$ es
\omterm{def:b3:topology:topology}{abierto}. \omterm{def:b3:topology:hausdorff}{De Hausdorff}: sea
$\bar x \neq \bar y$; elíjanse representantes con
$\delta = \min_{n}\abs{x - y - n} > 0$; las imágenes de los intervalos
de radio $\delta/3$ alrededor de $x$ y $y$ son abiertas (por ser $\pi$
abierta), contienen $\bar x, \bar y$ y son disjuntas (dos puntos de las
preimágenes distarían $< 2\delta/3$ módulo $\Z$).
\omterm{def:b3:topology:compact}{Compacto}: $\R/\Z = \pi([0,1])$, imagen
\omterm{def:b3:topology:continuity}{continua} de un
\omterm{def:b3:topology:compact}{espacio compacto}
(el~\cref{thm:b3:topology:compactprops}(2)). Así pues, $\bar\varphi$ es
una biyección \omterm{def:b3:topology:continuity}{continua} de un
\omterm{def:b3:topology:compact}{espacio compacto} en el
\omterm{def:b3:topology:hausdorff}{de Hausdorff} $S^1$: un
\omterm{def:b3:topology:continuity}{homeomorfismo}
(\cref{cor:b3:topology:compacthomeo}).

Para (iii): la compuesta $[0,1] \hookrightarrow \R
\xrightarrow{\pi} \R/\Z$ es
\omterm{def:b3:topology:continuity}{continua}, sobreyectiva e identifica
exactamente $0 \sim 1$: induce una biyección
\omterm{def:b3:topology:continuity}{continua} $[0,1]/(0\sim1) \to \R/\Z$
desde un \omterm{def:b3:topology:compact}{espacio compacto} (imagen
\omterm{def:b3:topology:continuity}{continua} de $[0,1]$ por la
proyección al cociente) a uno \omterm{def:b3:topology:hausdorff}{de
Hausdorff}: un \omterm{def:b3:topology:continuity}{homeomorfismo}.
Componiendo: los tres espacios son
\omterm{def:b3:topology:continuity}{homeomorfos}.
\end{solution}

\begin{solution}{exo:b3:topology:6}
(a) Un recubrimiento de $K_1 \cup \dots \cup K_m$ se restringe a un
recubrimiento de cada $K_j$: bastan un número finito de \omterm{def:b3:topology:topology}{abiertos} por
pieza. Una intersección $\bigcap K_i$ es cerrada en el
\omterm{def:b3:topology:compact}{compacto} $K_{i_0}$ (los
\omterm{def:b3:topology:compact}{compactos} son cerrados en el espacio
métrico ambiente), luego compacta.

(b) $x \mapsto d(x, F)$ es \omterm{def:b3:topology:continuity}{continua}
($1$-lipschitziana) y estrictamente positiva en $K$ ($d(x, F) = 0$
significaría $x \in \bar F =
F$); sobre el \omterm{def:b3:topology:compact}{compacto} $K$ alcanza un
mínimo $m > 0$: $d(K, F)
\geq m$. Contraejemplo sin \omterm{def:b3:topology:compact}{compacidad}: $F_1 = \{n : n \geq
2\}$ y $F_2 = \{n + \frac1n : n \geq 2\}$ son cerrados, disjuntos y
están a distancia $\inf_n \frac1n = 0$.

(c) Si $X$ es \omterm{def:b3:topology:compact}{compacto}, toda $f$
\omterm{def:b3:topology:continuity}{continua} está acotada
(el~\cref{thm:b3:topology:compactprops}(2)). Recíprocamente, si $X$ no
es \omterm{def:b3:topology:compact}{compacto}, tómese $(x_n)$ sin
subsucesión convergente (el~\cref{thm:b3:topology:metriccompact});
pasando a una subsucesión podemos suponer los $x_n$ distintos dos a dos,
y ningún punto de $X$ es punto de acumulación de la sucesión, de modo
que
\[
r_n = \min\Bigl(2^{-n},\ \tfrac13\,d\bigl(x_n, \{x_m : m \neq
n\}\bigr)\Bigr) > 0,
\]
y las bolas $B(x_n, r_n)$ son disjuntas dos a dos (un punto común de la
$n$-ésima y la $m$-ésima daría $d(x_n, x_m) < r_n +
r_m \leq \frac23 d(x_n, x_m)$). Defínase
\[
f(x) = \sum_{n} n\,\max\Bigl(0,\ 1 - \frac{d(x,
x_n)}{r_n}\Bigr):
\]
en cada $x$ hay a lo sumo un sumando no nulo, y $f(x_n) = n$.
\omterm{def:b3:topology:continuity}{Continuidad} en $x$: sea
$y_k \to x$; cada valor $f(y_k)$ es $0$ o el valor
$\mathrm{bump}_{n_k}(y_k)$ de una sola joroba. Si algún índice $n$
aparece infinitas veces, a lo largo de esa subsucesión
$f(y_k) = \mathrm{bump}_n(y_k) \to
\mathrm{bump}_n(x)$, que es igual a $f(x)$ (si $x \in B(x_n,
r_n)$, toda la cola está en esa bola abierta y no aparece ningún otro
índice; si $x \notin B(x_n, r_n)$, entonces
$\mathrm{bump}_n(x) = 0 = f(x)$, pues $x$, adherente a la $n$-ésima
bola, no está en ninguna otra bola abierta). Si $n_k \to \infty$:
$d(y_k, x_{n_k}) < r_{n_k} \leq 2^{-n_k} \to 0$, luego $x_{n_k}
\to x$, lo que haría de $x$ un punto de acumulación de $(x_n)$ ---
excluido; así que este caso solo afecta a un número finito de $k$. En
todos los casos $f(y_k) \to f(x)$: $f$ es
\omterm{def:b3:topology:continuity}{continua} y no acotada.
\end{solution}

\begin{solution}{exo:b3:topology:7}
Supongamos que ningún $\delta$ sirve: para cada $n$ hay un $A_n$ de
diámetro $< 1/n$ contenido en ningún $U_i$; tómese $x_n \in
A_n$. Por \omterm{def:b3:topology:compact}{compacidad} (el~\cref{thm:b3:topology:metriccompact}), una
subsucesión $x_{n_k} \to x$; tómense $i$ y $r > 0$ con $B(x,
r) \subseteq U_i$. Para $k$ grande: $d(x_{n_k}, x) < r/2$ y
$1/n_k < r/2$, luego $A_{n_k} \subseteq B(x_{n_k}, 1/n_k)
\subseteq B(x, r) \subseteq U_i$ --- contradicción. Heine: dado
$\varepsilon$, recúbrase $Y$ por bolas $B(y, \varepsilon/2)$; las
preimágenes forman un recubrimiento \omterm{def:b3:topology:topology}{abierto} de $X$; sea $\delta$ un
número de Lebesgue: si $d(x, x') < \delta$, el par $\{x, x'\}$ tiene
diámetro $< \delta$, está en una sola preimagen y $d(f(x),
f(x')) < \varepsilon$.
\end{solution}

\begin{solution}{exo:b3:topology:8}
(a) $GL_n(\R) = \det^{-1}(\R\setminus\{0\})$ es \omterm{def:b3:topology:topology}{abierto} ($\det$ es
polinómico, luego \omterm{def:b3:topology:continuity}{continuo}).
Disconexo: $\det$ lo aplica sobre $\R\setminus\{0\}$, y un
\omterm{def:b3:topology:connected}{espacio conexo} tiene imágenes
\omterm{def:b3:topology:continuity}{continuas}
\omterm{def:b3:topology:connected}{conexas}
(el~\cref{thm:b3:topology:connectedbasics}(2)); y $\R\setminus\{0\}$ no
es un intervalo. $GL_n(\C)$: para $A, B$ invertible,
$p(\lambda) = \det\bigl((1-\lambda)A + \lambda
B\bigr)$ es un polinomio en $\lambda$ con $p(0), p(1) \neq 0$, luego no
idénticamente nulo: tiene un número finito de raíces en $\C$. El plano
menos un número finito de puntos es
\omterm{def:b3:topology:pathconnected}{conexo por caminos} (esquívense
los puntos rodeándolos), de modo que hay un
\omterm{def:b3:topology:pathconnected}{camino} $\lambda(t)$ de $0$ a $1$
con $p(\lambda(t)) \neq 0$: $t \mapsto
(1-\lambda(t))A + \lambda(t)B$ es un
\omterm{def:b3:topology:pathconnected}{camino} en $GL_n(\C)$.

(b) $\det(A + \varepsilon I) = \chi_A(-\varepsilon)\cdot(\pm1)$
se anula para a lo sumo $n$ valores de $\varepsilon$: hay matrices
invertibles $A + \varepsilon I \to A$ con $\varepsilon \to 0$: densidad.
\end{solution}

\begin{solution}{exo:b3:topology:9}
(a) Sea $A \subseteq \Q$ que contiene $p < q$ y tómese un irracional
$\alpha \in (p, q)$: $A = (A \cap (-\infty, \alpha)) \sqcup
(A\cap(\alpha, +\infty))$ separa $A$ en dos piezas relativamente
abiertas no vacías: $A$ es disconexo. Las componentes se reducen a
puntos. La \omterm{def:b3:topology:compact}{compacidad} local falla: un
\omterm{def:b3:topology:topology}{entorno}
\omterm{def:b3:topology:compact}{compacto} $V$ de $0$ en $\Q$ contendría
$[-r, r]\cap\Q$ para cierto $r > 0$, que es cerrado en $V$ y, por tanto,
\omterm{def:b3:topology:compact}{compacto}; pero una sucesión de
racionales de $[-r,r]$ que converja (en $\R$) a un irracional no tiene
ninguna subsucesión convergente \emph{en $\Q$}: contradicción con
\cref{thm:b3:topology:metriccompact}.

(b) Sea $\gamma = (\gamma_1, \gamma_2) \colon [0,1] \to S$ un
\omterm{def:b3:topology:pathconnected}{camino} con $\gamma(0) = (0,0)$ y
$\gamma(1) \in \Gamma$. El conjunto $\{t : \gamma_1(t) = 0\}$ es cerrado
y no contiene $1$; sea $t^*$ su supremo, de modo que $\gamma_1(t^*) = 0$
y $\gamma_1 > 0$ sobre $(t^*, 1]$. Por la
\omterm{def:b3:topology:continuity}{continuidad} de $\gamma_2$ en $t^*$,
elíjase $\delta$ con $\abs{\gamma_2(t) -
\gamma_2(t^*)} < \tfrac12$ para $t \in [t^*, t^* + \delta]$. Fijemos
$t_0 = t^* + \delta$ (podemos suponer $\leq 1$): $\gamma_1(t_0) > 0$, y
$\gamma_1$ toma todos los valores de $(0, \gamma_1(t_0)]$ sobre
$[t^*, t_0]$ (teorema del valor intermedio,
el~\cref{thm:b3:topology:connectedbasics}). Elíjase $k$ grande con
$a_k = \frac1{\pi/2 + 2k\pi} < \gamma_1(t_0)$ y $b_k =
\frac1{3\pi/2 + 2k\pi} < \gamma_1(t_0)$: existen $s, s' \in
(t^*, t_0]$ con $\gamma_1(s) = a_k$, $\gamma_1(s') = b_k$; como los
puntos están sobre la gráfica, $\gamma_2(s) = \sin(1/a_k) =
1$ y $\gamma_2(s') = -1$. Ambos no pueden distar menos de $\frac12$ de
$\gamma_2(t^*)$: contradicción. $S$ es
\omterm{def:b3:topology:connected}{conexo} pero no
\omterm{def:b3:topology:pathconnected}{conexo por caminos}.
\end{solution}

\begin{solution}{exo:b3:topology:10}
(a) $C_n$ consta exactamente de los $x$ que admiten un desarrollo
ternario con dígitos en $\{0, 2\}$ hasta el rango $n$ (inducción:
suprimir los tercios centrales elimina el primer dígito $1$, etc.; los
extremos tienen dos desarrollos, uno de ellos sin $1$), de modo que
$C = \bigcap C_n$ es el conjunto de las sumas $\sum a_n3^{-n}$ con
$a_n \in \{0,2\}$. \omterm{def:b3:topology:compact}{Compacto}: cada
$C_n$ es unión finita de intervalos cerrados, y $C$ es cerrado en
$[0,1]$. \omterm{def:b3:topology:interior}{Interior} vacío:
$C \subseteq C_n$, unión de intervalos de longitud $3^{-n}$; un
intervalo \omterm{def:b3:topology:interior}{interior} de longitud
$\ell > 0$ tendría que caber en uno de ellos para todo $n$.
\omterm{prop:b3:galois:perfect}{Perfecto}: dados $x \in C$ y $m$,
cámbiese el dígito $a_m$: el nuevo punto está en $C$, es distinto y
dista menos de $2\cdot3^{-m}$ de $x$.

(b) Si $x \neq y \in C$, sus desarrollos difieren por primera vez en
cierto rango $k$; entre ellos hay un tercio central suprimido (el hueco
de rango $k$ que separa el dígito $0$ del dígito $2$), lo que
proporciona un punto $c \notin C$, $x < c < y$ (digamos): $C =
(C \cap (-\infty, c)) \sqcup (C \cap (c, \infty))$ desconecta cualquier
subconjunto que contenga ambos puntos. Las componentes se reducen a
puntos.

(c) La aplicación de dígitos $\beta\colon x \mapsto (a_n(x))$ es una
biyección sobre $\{0,2\}^\N$ (existencia y unicidad de los desarrollos
en \omterm{def:b3:topology:basis}{base} $\{0,2\}$: sucesiones de dígitos distintas dan puntos a
distancia $\geq 3^{-k} > 0$ del rango en que difieren por primera vez,
como en el~\cref{pb:b3:topology:1}, pregunta 1). Es
\omterm{def:b3:topology:continuity}{continua}: $\abs{x - y} < 3^{-m}$
obliga a que coincidan los $m$ primeros dígitos (mismo cálculo), de modo
que $\beta$ lleva bolas pequeñas dentro de cajas de la
\omterm{def:b3:topology:basis}{base}. Una biyección
\omterm{def:b3:topology:continuity}{continua} del
\omterm{def:b3:topology:compact}{compacto} $C$ al producto
\omterm{def:b3:topology:hausdorff}{de Hausdorff} es un
\omterm{def:b3:topology:continuity}{homeomorfismo}
(el~\cref{cor:b3:topology:compacthomeo}; el producto es
\omterm{def:b3:topology:hausdorff}{de Hausdorff}: sepárese en una
coordenada en la que difieran). No numerabilidad: diagonal de Cantor
sobre $\{0,2\}^\N$. Por último, $\{0,2\}^\N \times
\{0,2\}^\N \cong \{0,2\}^\N$ entrelazando dígitos (un
\omterm{def:b3:topology:continuity}{homeomorfismo}:
\omterm{def:b3:topology:continuity}{continuidad} componente a componente
en ambos sentidos), luego $C
\times C \cong C$.
\end{solution}

\begin{solution}{exo:b3:topology:11}
Para $y \in \R^n$, póngase
\[
\tau(y) = \Bigl(\frac{2y}{\norm y^2 + 1},\ \frac{\norm y^2 -
1}{\norm y^2 + 1}\Bigr) \in \R^n\times\R = \R^{n+1}.
\]
Se comprueba que $\norm{\tau(y)} = 1$, $\tau(y) \neq N$, y
$\sigma(\tau(y)) = y$, $\tau(\sigma(x)) = x$: $\sigma$ y $\tau$ son
mutuamente inversas y ambas
\omterm{def:b3:topology:continuity}{continuas} (fórmulas racionales con
denominadores que no se anulan): $S^n\setminus\{N\} \cong
\R^n$. Extiéndase a $\hat\sigma \colon S^n \to \widehat{\R^n}$ mediante
$N \mapsto \infty$: es una biyección,
\omterm{def:b3:topology:continuity}{continua} en todo punto de
$S^n\setminus\{N\}$ y en $N$: un
\omterm{def:b3:topology:topology}{entorno} de la
\omterm{def:b3:topology:basis}{base} de $\infty$ es
$\widehat{\R^n}\setminus K$ con $K$
\omterm{def:b3:topology:compact}{compacto}, $K
\subseteq \bar B(0, R)$; como $\norm{\sigma(x)}^2 = \frac{1 +
x_{n+1}}{1 - x_{n+1}} \to \infty$ cuando $x_{n+1} \to 1$, el conjunto
$\{x \in S^n : x_{n+1} > 1 - \eta\}$ se aplica fuera de $\bar B(0,R)$
para $\eta$ pequeño: \omterm{def:b3:topology:continuity}{continuidad}.
Una biyección \omterm{def:b3:topology:continuity}{continua} del
\omterm{def:b3:topology:compact}{compacto} $S^n$ al
\omterm{def:b3:topology:hausdorff}{de Hausdorff} $\widehat{\R^n}$ es un
\omterm{def:b3:topology:continuity}{homeomorfismo}. Al suprimir otro
punto $p$: una rotación de $S^n$ lleva $p$ a $N$ (las rotaciones son
\omterm{def:b3:topology:continuity}{homeomorfismos}), lo que reduce al
caso ya calculado: $S^n \setminus \{p\} \cong \R^n$.
\end{solution}

\begin{solution}{exo:b3:topology:12}
(a) $T$ está acotado y es cerrado: un límite de puntos de $T$ con
abscisas $\to x_0 > 0$ permanece sobre la gráfica (localmente cerrada)
por la \omterm{def:b3:topology:continuity}{continuidad} de $\sin\frac1x$
en $\intoc01$; un límite con abscisas $\to 0$ tiene ordenada en
$\intcc{-1}1$ y, por tanto, está en el segmento.
\omterm{def:b3:topology:compact}{Compacto} por Heine--Borel
(el~\cref{cor:b3:topology:heineborel}).
\omterm{def:b3:topology:interior}{Clausura} de la gráfica $\Gamma$: todo
punto $(0, y)$ con $y \in \intcc{-1}1$ es límite de puntos de la gráfica
--- resuélvase $\sin\frac1x = y$ cerca de $0$ (la función barre
$\intcc{-1}1$ en cada intervalo
$[\frac1{(k+1)\pi}, \frac1{k\pi}]$): $\bar\Gamma = T$.

(b) $\Gamma$ es la imagen \omterm{def:b3:topology:continuity}{continua}
del \omterm{def:b3:topology:connected}{conexo} $\intoc01$ por
$x \mapsto (x, \sin\frac1x)$: es
\omterm{def:b3:topology:connected}{conexa}; y la
\omterm{def:b3:topology:interior}{clausura} de un
\omterm{def:b3:topology:connected}{conexo} es
\omterm{def:b3:topology:connected}{conexa}
(\cref{thm:b3:topology:connectedbasics}): $T = \bar\Gamma$
es \omterm{def:b3:topology:connected}{conexo}.

(c) Supongamos que $\gamma\colon\intcc01\to T$ es
\omterm{def:b3:topology:continuity}{continua} con $\gamma(0) = (0,0)$,
$\gamma(1) = (1, \sin1)$, y sea $t^*
= \sup\{t : \gamma_1(t) = 0\}$: por
\omterm{def:b3:topology:continuity}{continuidad}, $\gamma_1(t^*)
= 0$ y $\gamma_1 > 0$ sobre $\intoc{t^*}1$. Para todo $\delta > 0$,
sobre el intervalo $\intoc{t^*}{t^*+\delta}$ la función $\gamma_1$ toma
todos los valores de cierto $\intoc0\eta$ (teorema del valor intermedio,
$\gamma_1(t^* + \delta) > 0$); en particular, contiene abscisas de la
forma $\frac1{\pi/2 + 2k\pi}$ y $\frac1{-\pi/2 + 2k\pi}$ con $k$
arbitrariamente grande, en las que $\gamma_2 = \pm1$. Así, en todo
\omterm{def:b3:topology:topology}{entorno} por la derecha de $t^*$,
$\gamma_2$ toma ambos valores $1$ y $-1$: $\gamma_2$ no tiene límite en
$t^{*+}$, en contra de la
\omterm{def:b3:topology:continuity}{continuidad}. No existe tal
\omterm{def:b3:topology:pathconnected}{camino}.

(d) $T$ es \omterm{def:b3:topology:connected}{conexo} pero no
\omterm{def:b3:topology:pathconnected}{conexo por caminos}: las dos
nociones difieren. Para $U \subseteq \R^n$ \omterm{def:b3:topology:topology}{abierto} y
\omterm{def:b3:topology:connected}{conexo} y $a \in U$, sea $V$ el
conjunto de los puntos de $U$ unibles a $a$ por un
\omterm{def:b3:topology:pathconnected}{camino} en $U$. $V$ es \omterm{def:b3:topology:topology}{abierto}:
alrededor de $x \in V$ hay una bola $B(x, r)
\subseteq U$ estrellada, y concatenando el
\omterm{def:b3:topology:pathconnected}{camino} hasta $x$ con un segmento
se alcanza todo punto de la bola. $V$ es cerrado en $U$: si
$x \in U \setminus V$, el mismo argumento de la bola muestra que
$B(x, r) \cap V = \varnothing$ (un punto de $V$ en la bola uniría $a$
con $x$). No vacío ($a \in
V$), \omterm{def:b3:topology:topology}{abierto} y cerrado en el \omterm{def:b3:topology:connected}{conexo}
$U$: $V = U$. La curva seno escapa a esto por ser cerrada de
\omterm{def:b3:topology:interior}{interior} vacío: su punto «malo»
$(0,0)$ no tiene ninguna bola dentro de $T$ con la que salvar las
oscilaciones.
\end{solution}

\begin{solution}{pb:b3:topology:1}
\textbf{1.} Supongamos que los desarrollos de $x, y \in C$ difieren por
primera vez en el rango $k \leq m$, digamos $b_k(x) = 1$, $b_k(y) = 0$.
Entonces
\[
x - y \;\geq\; \frac{2}{3^k} - \sum_{j > k}\frac{2}{3^j}
= \frac{2}{3^k} - \frac{1}{3^k} = \frac1{3^k} \geq \frac1{3^m}:
\]
contrarrecíproco: $\abs{x - y} < 3^{-m}$ obliga a que coincidan hasta el
rango $m$.

\textbf{2.} Por la pregunta 1, $b_n$ es constante sobre $C \cap (x -
3^{-n}, x + 3^{-n})$: es localmente constante y, por tanto,
\omterm{def:b3:topology:continuity}{continua}. La aplicación
$x \mapsto (b_n(x))_n \in \{0,1\}^\N$ es
\omterm{def:b3:topology:continuity}{continua} (componente a componente,
la~\cref{prop:b3:topology:universal}(a)) y biyectiva (unicidad de los
desarrollos de dígitos $\{0,2\}$); desde el
\omterm{def:b3:topology:compact}{compacto} $C$ a un
\omterm{def:b3:topology:hausdorff}{espacio de Hausdorff}, es un
\omterm{def:b3:topology:continuity}{homeomorfismo}
(\cref{cor:b3:topology:compacthomeo}).

\textbf{3.} \omterm{def:b3:topology:continuity}{Continuidad}: si
$\abs{x - y} < 3^{-m}$, los $m$ primeros dígitos coinciden, luego
$\abs{g(x) - g(y)} \leq \sum_{n > m}2^{-n}
= 2^{-m}$. Sobreyectividad: todo $t \in [0,1]$ tiene un desarrollo
binario $t = \sum b_n2^{-n}$, y $t = g\bigl(\sum
2b_n3^{-n}\bigr)$. No inyectiva: $g$ identifica los dos puntos de Cantor
de dígitos $(0,1,1,1,\dots)$ y $(1,0,0,\dots)$ --- ambos van a $\frac12$
(la ambigüedad diádica $0.0111\ldots_2 =
0.1000\ldots_2$).

\textbf{4.} Cada componente de $\Phi$ es
\omterm{def:b3:topology:continuity}{continua} por la misma estimación
(sus dígitos son una subsucesión de los $b_n$). Sobreyectividad: dado
$(s, t) \in [0,1]^2$, elíjanse dígitos binarios $(c_n)$ de $s$ y $(d_n)$
de $t$, y entrelácense: el punto $x
\in C$ con $b_{2n-1}(x) = c_n$, $b_{2n}(x) = d_n$ cumple
$\Phi(x) = (s,t)$.

\textbf{5.} Los huecos son disjuntos dos a dos y tienen los extremos en
$C$, de modo que la fórmula define $f$ sin ambigüedad sobre $[0,1]$; en
los extremos de un hueco, la fórmula afín devuelve $\Phi(u)$, $\Phi(v)$:
coherente con $f = \Phi$ sobre $C$. Sobreyectividad: ya lo era
$f(C) = \Phi(C) = [0,1]^2$.

\textbf{6.} En $t_0 \notin C$: $t_0$ está en un hueco \omterm{def:b3:topology:topology}{abierto} sobre el
que $f$ es afín: \omterm{def:b3:topology:continuity}{continua}. En
$x \in C$: registremos primero dos estimaciones.

\emph{(i) Sobre $C$}: si $\abs{y - x} \leq 3^{-M}$, $y \in C$, entonces
los dígitos coinciden hasta $M$, de modo que cada componente de $\Phi(y)
- \Phi(x)$ es a lo sumo $\sum_{n > \lfloor M/2\rfloor} 2^{-n} =
2^{-\lfloor M/2\rfloor}$.

\emph{(ii) A través de un hueco}: un hueco $(u,v)$ suprimido en la etapa
$k$ tiene longitud $3^{-k}$, y los dígitos de sus extremos coinciden
hasta el rango $k - 1$ (difieren a partir del rango $k$), de modo que
$\norm{\Phi(v) - \Phi(u)}_\infty \leq
2^{-\lfloor (k-1)/2\rfloor}$.

Sean ahora $\abs{t - x} < 3^{-M}$ con $t$ en un hueco $(u, v)$ de etapa
$k$, y sea $u$ el extremo del lado de $x$, de forma que
$\abs{u - x} < 3^{-M}$ y $0 < t - u < 3^{-M}$. Entonces
\[
\norm{f(t) - \Phi(x)}_\infty
\leq \norm{\Phi(u) - \Phi(x)}_\infty
+ \frac{t - u}{v - u}\,\norm{\Phi(v) - \Phi(u)}_\infty
\leq 2^{-\lfloor M/2\rfloor} + \min\bigl(1,\ 3^{k - M}\bigr)\,
2^{-\lfloor(k-1)/2\rfloor} .
\]
Para $k \geq M$, el último término es $\leq 2^{-\lfloor(M-1)/2
\rfloor}$; para $k < M$, como $2^{-\lfloor(k-1)/2\rfloor} \leq
\sqrt2\,\cdot 2^{-k/2}$,
\[
3^{k-M}\,2^{-\lfloor(k-1)/2\rfloor}
\leq \sqrt2\,\Bigl(\frac{3}{\sqrt2}\Bigr)^{k} 3^{-M}
\leq \sqrt2\,\Bigl(\frac{3}{\sqrt2}\Bigr)^{M} 3^{-M}
= \sqrt2\; 2^{-M/2}
\]
(la cantidad intermedia crece con $k$). En todos los casos,
$\norm{f(t) - \Phi(x)}_\infty \leq C\,2^{-M/2}$
con una constante absoluta $C$: haciendo $M \to \infty$ se demuestra la
\omterm{def:b3:topology:continuity}{continuidad} en $x$ (los puntos
$t \in C$ los cubre (i)). Por tanto, $f$ es una sobreyección
\omterm{def:b3:topology:continuity}{continua} $[0,1] \to [0,1]^2$ --- de
hecho, las estimaciones muestran que es de tipo Hölder con exponente
$\log 2/(2\log
3)$, pero la \omterm{def:b3:topology:continuity}{continuidad} es todo lo
que hemos afirmado.

\textbf{7.} Para $[0,1]^n$: sepárense los dígitos de $x \in C$ en $n$
subsucesiones entrelazadas y repítanse las preguntas 4--6 literalmente.
Para $\R \to \R^2$: sea $h \colon [0,1] \to [-1,1]^2$ una sobreyección
\omterm{def:b3:topology:continuity}{continua} (reescálese $f$). Defínase
$F$ sobre $[m, m
+ \frac12]$ ($m \in \Z$) como una copia de $h$ reescalada para llenar
$[-(m+1), m+1]^2$ si $m \geq 0$ (y simétricamente para $m <
0$), y sobre $[m + \frac12, m + 1]$ como el segmento afín que une los
valores en los extremos: $F$ es
\omterm{def:b3:topology:continuity}{continua} (pegado sobre cerrados,
la~\cref{prop:b3:topology:gluing}(b)) y su imagen contiene
$\bigcup_m [-(m+1), m+1]^2 = \R^2$.

\textbf{8.} $[0,1]$ es \omterm{def:b3:topology:compact}{compacto} y
$[0,1]^2$ es \omterm{def:b3:topology:hausdorff}{de Hausdorff}: una
inyección \omterm{def:b3:topology:continuity}{continua} es un
\omterm{def:b3:topology:continuity}{homeomorfismo} sobre su imagen
(el~\cref{cor:b3:topology:compacthomeo} aplicado a la correstricción).

\textbf{9.} Suprimamos $\frac12$: $[0,1]\setminus\{\frac12\}$ es
disconexo. Si $h \colon [0,1] \to [0,1]^2$ fuera un
\omterm{def:b3:topology:continuity}{homeomorfismo},
$[0,1]^2\setminus\{h(\frac12)\}$ también sería disconexo (las imágenes
\omterm{def:b3:topology:continuity}{homeomorfas} de espacios disconexos
son disconexas); pero el cuadrado menos un punto es
\omterm{def:b3:topology:pathconnected}{conexo por caminos}: únanse dos
puntos por un \omterm{def:b3:topology:pathconnected}{camino} de dos
segmentos que evite el punto suprimido. Contradicción:
$[0,1] \not\cong [0,1]^2$.

\textbf{10.} Por la pregunta 8, una sobreyección
\omterm{def:b3:topology:continuity}{continua} inyectiva
$[0,1] \to [0,1]^2$ sería un
\omterm{def:b3:topology:continuity}{homeomorfismo}, en contra de la
pregunta 9. La \omterm{def:b3:topology:compact}{compacidad} intervino exactamente en
el~\cref{cor:b3:topology:compacthomeo}: es lo que hace
\omterm{def:b3:topology:continuity}{continua} la inversa de la biyección
\omterm{def:b3:topology:continuity}{continua} (las imágenes de cerrados
son compactas y, por tanto, cerradas).
Sin \omterm{def:b3:topology:compact}{compacidad} la conclusión falla realmente: $[0, 2\pi) \to S^1$ es una
biyección \omterm{def:b3:topology:continuity}{continua} que no es un
\omterm{def:b3:topology:continuity}{homeomorfismo}.

\textbf{11.} Aquí queda demostrado: \emph{existe una sobreyección
\omterm{def:b3:topology:continuity}{continua}, pero ninguna biyección
\omterm{def:b3:topology:continuity}{continua}, de $[0,1]$ sobre
$[0,1]^2$; en particular, $[0,1] \not\cong [0,1]^2$.} Así pues, la
«dimensión» no se conserva por sobreyecciones
\omterm{def:b3:topology:continuity}{continuas} --- ni el cardinal ni
siquiera la \omterm{def:b3:topology:continuity}{continuidad} pueden
verla ---, pero \emph{sí} es un invariante
\omterm{def:b3:topology:topology}{topológico} al nivel de los
\omterm{def:b3:topology:continuity}{homeomorfismos}, al menos para las
dimensiones $1$ y $2$, donde basta la \omterm{def:b3:topology:connected}{conexión} tras suprimir un punto.
El enunciado general --- la \emph{invariancia del dominio}:
$\R^m \cong \R^n$ implica $m =
n$, y una inyección \omterm{def:b3:topology:continuity}{continua}
$\R^n \to \R^n$ es abierta --- es cierto, pero exige
\omterm{def:b3:topology:topology}{topología}
\omterm{def:b3:galois:algebraic}{algebraica} (homología), fuera de este
curso.

\textbf{12.} En \omterm{def:b3:topology:basis}{base} ternaria, $\frac C2 = \{\sum_na_n3^{-n} : a_n
\in \{0, 1\}\}$ (dividir por dos los dígitos $\{0,2\}$ da los dígitos
$\{0,1\}$). Para $t \in [0,1]$, tómese cualquier desarrollo ternario
$t =
\sum_nd_n3^{-n}$, $d_n \in \{0,1,2\}$, y sepárese cada dígito como
$d_n = a_n + b_n$ con $a_n, b_n \in \{0,1\}$ ($0 = 0{+}0$, $1 = 1{+}0$,
$2 = 1{+}1$): entonces $t \in \frac C2 + \frac C2$. El meollo es que
cada dígito se separa \emph{dentro} de $\{0,1\}$, de modo que ningún
acarreo se propaga y los dígitos pueden tratarse por separado. La
inclusión recíproca es clara (las sumas de dígitos siguen siendo
$\leq 2$). Reescalando por $2$: $C + C = [0,2]$.

\textbf{13.} La aplicación $(x, y) \mapsto x + y$ lleva $C \times C$
sobre $C + C = [0,2]$: el polvo, que no contiene ningún cuadrado
(\omterm{def:b3:topology:interior}{interior} vacío:
el~\cref{exo:b3:topology:10}), se proyecta a lo largo de la antidiagonal
sobre un segmento completo. Autosemejanza: el primer dígito ternario de
un punto de Cantor es $0$ o $2$, y al quitarlo queda
$C = \frac C3 \cup \bigl(\frac C3 +
\frac23\bigr)$ --- la ecuación de punto fijo que genera toda imagen de
$C$.

\textbf{14.} A través de $C \cong \{0,1\}^\N$ (pregunta 2), el
entrelazado $(\varepsilon, \varepsilon') \mapsto
(\varepsilon_1, \varepsilon_1', \varepsilon_2,
\varepsilon_2', \dots)$ es una biyección
$\{0,1\}^\N\times\{0,1\}^\N \to \{0,1\}^\N$,
\omterm{def:b3:topology:continuity}{continua} en ambos sentidos (cada
coordenada de salida depende de una sola coordenada de entrada;
\omterm{def:b3:topology:constructions}{topología producto}). Por tanto,
$C\times C \cong C$ y, por inducción, $C^n \cong C$. Los espacios
familiares no comparten esto: $\R^2 \not\cong \R$ y
$[0,1]^2 \not\cong [0,1]$ (pregunta 11); los productos infinitos sí:
$([0,1]^\N)^2 \cong
[0,1]^\N$ por el mismo entrelazado.

\textbf{15.} $\sum_{k\geq1}2\cdot3^{-2k} = 2\cdot\frac{1/9}{1
- 1/9} = \frac14$: el desarrollo ternario de $\frac14$ es
$0.\overline{02}$, con dígitos en $\{0,2\}$, luego $\frac14 \in
C$; y como no es un desarrollo finito ni tiene cola constante $2$, no es
extremo de ningún tercio central suprimido. Los extremos forman un
conjunto numerable (dos por intervalo suprimido, y una cantidad
numerable de intervalos), mientras que $C \cong \{0,1\}^\N$ es no
numerable (argumento diagonal, o pregunta 19): casi todo punto de Cantor
es, como $\frac14$, invisible en la imagen habitual de los extremos.

\textbf{16.} $g(x) = g(y)$ significa que las sucesiones binarias
$(b_n(x))$, $(b_n(y))$ representan el mismo real. Las sucesiones
binarias distintas que representan el mismo número son exactamente las
de la ambigüedad diádica $0.w01^\infty = 0.w10^\infty$: hay a lo sumo
dos preimágenes, y exactamente dos precisamente en los racionales
diádicos de $(0,1)$. Así pues, $g$ colapsa $C$ sobre $[0,1]$ pegando una
cantidad numerable de parejas --- el esqueleto combinatorio de la
escalera de Cantor del~\cref{ch:b3:measure}.

\textbf{17.} $x$ no es aislado en $P$, de modo que $B(x, r/2)$ contiene
algún $y_0 \in P \setminus \{x\}$; $y_0$ tampoco es aislado, de modo que
$B(x, r/2)$ contiene un segundo punto $y_1
\in P$. Cualquier $\rho \leq \min\bigl(\frac{d(y_0,y_1)}3,
\frac r4\bigr)$ hace $\bar B(y_0, \rho)$ y $\bar B(y_1,
\rho)$ disjuntas y contenidas en $B(x, r)$. Ambas están centradas en
puntos de $P$, donde puede repetirse la misma extracción de dos puntos:
la perfección es la reserva inagotable de puntos próximos que mantiene
viva la recursión para siempre.

\textbf{18.} Constrúyanse los $F_w$ por inducción sobre la longitud de
la palabra: $F_\varnothing = P$, y dentro de cada bola $B(y_w, \rho_w)$
la pregunta 17 proporciona dos bolas cerradas disjuntas de radio $\leq
\min(\rho_w/2,\ 2^{-\abs w-1}\operatorname{diam}P)$ centradas en puntos
de $P$; póngase $F_{wi} = \bar B(y_{wi}, \rho_{wi})
\cap P$: es no vacío (su \omterm{ex:b3:groups:actions}{centro}) y
\omterm{def:b3:topology:compact}{compacto}. Para $\varepsilon$ infinita,
los \omterm{def:b3:topology:compact}{compactos} encajados no vacíos
$F_{\varepsilon_1\cdots\varepsilon_m}$ tienen intersección no vacía
(propiedad de intersección finita en el
\omterm{def:b3:topology:compact}{compacto} $P$) y de diámetro $0$: un
único punto $h(\varepsilon)$.

\textbf{19.} Dos palabras que difieren por primera vez en el rango $m$
envían sus imágenes a los dos conjuntos \emph{disjuntos}
$F_{w0}, F_{w1}$ (prefijo común $w$): $h$ es inyectiva. Si dos palabras
coinciden hasta el rango $m$, ambas imágenes están en un conjunto de
diámetro $\leq
2^{-m}\operatorname{diam}P$: $h$ es
\omterm{def:b3:topology:continuity}{continua}. Así, el conjunto no
numerable $\{0,1\}^\N$ se inyecta en $P$: todo espacio métrico
\omterm{def:b3:topology:compact}{compacto}
\omterm{prop:b3:galois:perfect}{perfecto} es no numerable --- $[0,1]$ y
$C$ entre ellos.

\textbf{20.} $h$ es una inyección
\omterm{def:b3:topology:continuity}{continua} del
\omterm{def:b3:topology:compact}{compacto} $\{0,1\}^\N$ en el
\omterm{def:b3:topology:hausdorff}{de Hausdorff} (métrico) $P$:
el~\cref{cor:b3:topology:compacthomeo} la mejora a un
\omterm{def:b3:topology:continuity}{homeomorfismo} sobre su imagen; y
$\{0,1\}^\N \cong C$ (pregunta 2). Todo espacio métrico
\omterm{def:b3:topology:compact}{compacto}
\omterm{prop:b3:galois:perfect}{perfecto} contiene una copia del
conjunto de Cantor: la perfección tiene un germen universal, y es el de
Cantor.

\textbf{21.} Un espacio métrico
\omterm{def:b3:topology:compact}{compacto} numerable no puede ser
\omterm{prop:b3:galois:perfect}{perfecto} (pregunta 19), de modo que
tiene un punto aislado. En $K =
\{0\}\cup\{\frac1n : n \geq 1\}$, todo punto $\frac1n$ es aislado y $0$
no lo es: los puntos aislados pueden incluso ser
\omterm{def:b3:topology:interior}{densos} sin que el espacio sea
discreto --- la \omterm{def:b3:topology:compact}{compacidad} mantiene su límite dentro.

\textbf{22.} Sea $P$ el conjunto de los puntos de condensación de $F$ y
sea $\mathcal B$ la familia numerable de los intervalos \omterm{def:b3:topology:topology}{abiertos} con
extremos racionales. Todo punto de $F$ que no sea de condensación está
en algún $I \in \mathcal B$ con $I \cap F$ numerable, de modo que $F
\setminus P$ está contenido en la unión de esas trazas numerables, en
cantidad numerable: es numerable. Como $F$ es no numerable,
$P \neq \varnothing$; $P$ es cerrado (si $x \notin P$, el intervalo
testigo $I$ no contiene ningún punto de condensación: $I \cap F$ es
numerable) y $P \subseteq F$ ($F$ es cerrado: un punto de condensación
es en particular adherente). $P$ es
\omterm{prop:b3:galois:perfect}{perfecto}: si fuera $I \cap P = \{x\}$
para algún intervalo $I$, entonces
$I \cap F \subseteq \{x\} \cup (F \setminus P)$ sería numerable, en
contra de $x \in P$. Y la construcción del árbol de las preguntas 17--19
funciona literalmente dentro de $P$: las piezas $\bar B(y, \rho) \cap P$
son compactas (cerradas y acotadas en
$\R$), todo \omterm{ex:b3:groups:actions}{centro} es no aislado en $P$,
y el argumento no usó nada más. Por tanto, $\abs P \geq
\abs{\{0,1\}^\N} = \mathfrak c$, y $\abs F \leq \mathfrak
c$ trivialmente: un cerrado no numerable $F \subseteq \R$ tiene cardinal
exactamente $\mathfrak c$. Los cerrados no pueden atestiguar un fallo de
la hipótesis del \omterm{def:b3:topology:continuity}{continuo}.

\textbf{23.} Sean $x \neq y$ en $C$ y elíjase $m \geq 0$ con
$3^{-(m+1)} \leq \abs{x - y} < 3^{-m}$. Por la pregunta 1, los dígitos
$b_n(x) = b_n(y)$ coinciden para $n \leq m$; la primera coordenada de
$\Phi$ usa $b_1, b_3, \dots$ y la segunda $b_2,
b_4, \dots$, de modo que en cada coordenada los dos puntos comparten al
menos $\lfloor m/2\rfloor$ dígitos binarios iniciales:
\[
\norm{\Phi(x) - \Phi(y)}_\infty
\leq \sum_{k > \lfloor m/2\rfloor}2^{-k}
= 2^{-\lfloor m/2\rfloor} \leq \sqrt2\cdot2^{-m/2}
= \sqrt2\,\bigl(3^{-m}\bigr)^\alpha
\leq \sqrt2\,\bigl(3\abs{x-y}\bigr)^\alpha,
\]
usando $2^{-m/2} = 3^{-m\alpha}$ (tómense logaritmos); y
$\sqrt2\cdot3^\alpha = \sqrt2\cdot\sqrt2 = 2$: la cota
$2\abs{x-y}^\alpha$ sobre $C$. Mismo hueco $(u, v)$: $f$ es afín allí,
luego, para $t, s \in [u, v]$,
\[
\norm{f(t) - f(s)} = \frac{\abs{t-s}}{v-u}\,
\norm{\Phi(v) - \Phi(u)}
\leq \frac{\abs{t-s}}{v-u}\,2(v-u)^\alpha
\leq 2\abs{t-s}^\alpha,
\]
puesto que $\frac{\abs{t-s}}{v-u} \leq \bigl(\frac{\abs{t-s}}
{v-u}\bigr)^\alpha$ cuando $\abs{t-s} \leq v - u$. Caso general $t <
s$: si $[t, s]$ corta a $C$, sean $u'$ y $v'$ el menor y el mayor punto
del \omterm{def:b3:topology:compact}{compacto} $C \cap [t, s]$; entonces
$t$ está en un hueco (o en un punto de $C$) cuyo extremo derecho es
$u'$, y $s$ en uno cuyo extremo izquierdo es $v'$, y
\[
\norm{f(t) - f(s)} \leq 2\abs{t - u'}^\alpha
+ 2\abs{v' - u'}^\alpha + 2\abs{s - v'}^\alpha
\leq 6\abs{t - s}^\alpha ;
\]
si $[t, s]$ no corta a $C$, se aplica el caso del mismo hueco. Así pues,
$f$ es $\alpha$-Hölder con $\alpha = \frac{\ln2}{2\ln3}$.

\textbf{24.} Supongamos $\norm{F(t) - F(s)} \leq
C\abs{t-s}^\alpha$ con $F$ sobreyectiva y $\alpha >
\frac12$. Córtese $[0,1]$ en $N$ intervalos $I_1, \dots, I_N$ de
longitud $\frac1N$: cada imagen $F(I_j)$ tiene diámetro $\leq
CN^{-\alpha}$. Dos puntos distintos de la retícula $G_M =
\bigl\{\bigl(\frac iM, \frac jM\bigr)\bigr\}_{0 \leq i, j
\leq M}$ distan $\geq \frac1M$, de modo que un conjunto de diámetro
$< \frac1M$ contiene a lo sumo uno de ellos. Tómese $M =
\lfloor N^\alpha/(2C)\rfloor$ (con $N$ grande): entonces
$CN^{-\alpha} \leq \frac1{2M} < \frac1M$, y la sobreyectividad mete cada
uno de los $(M+1)^2$ puntos de la retícula en algún $F(I_j)$, de donde
\[
\frac{N^{2\alpha}}{4C^2} \leq (M + 1)^2 \leq N,
\]
y $N^{2\alpha - 1} \leq 4C^2$ falla para $N$ grande cuando
$2\alpha > 1$: contradicción. Nuestra curva, con $\alpha \approx
0.3155$, queda por debajo del muro, como debía; la curva de Hilbert
(admitido) es $\frac12$-Hölder, de modo que el exponente crítico
$\frac12$ sí se alcanza --- la regularidad Hölder, a diferencia de la
inyectividad, es cuestión de grado, y $\frac12$ es exactamente la
frontera que la dimensión $2$ impone a una aplicación procedente de la
dimensión $1$.

\textbf{25.} El $m$-ésimo racimo $\{\frac1m + \frac1n : n >
m^2\}$ está en $\bigl(\frac1m, \frac1m + \frac1{m^2}\bigr]
\subseteq \bigl(\frac1m, \frac1{m-1}\bigr)$ para $m \geq 2$, porque
$\frac1{m-1} - \frac1m = \frac1{m(m-1)} >
\frac1{m^2}$ (y el primer racimo está en $(1, \frac32]$): los racimos
ocupan intervalos disjuntos. Puntos de acumulación de $K$: dentro del
$m$-ésimo intervalo, solo $\frac1m$ (el racimo converge a él y es
discreto en sí mismo); globalmente, $0$ (todo
\omterm{def:b3:topology:topology}{entorno} de $0$ contiene racimos
enteros). Luego $K'
= \{0\} \cup \{\frac1m\}$, después $K'' = \{0\}$ (cada $\frac1m$ es
aislado en $K'$) y $K''' = \varnothing$; $K$ contiene sus puntos de
acumulación, luego es cerrado, acotado, numerable y
\omterm{def:b3:topology:compact}{compacto}, y sus puntos aislados ---
los puntos de los racimos $\frac1m +
\frac1n$ --- son \omterm{def:b3:topology:interior}{densos} en $K$: todo
elemento de $K'$ es límite suyo, tal como predice el mecanismo de la
pregunta 21. Inducción: $K_1 = \{0\}\cup\{\frac1n\}$ cumple
$K_1^{(1)} = \{0\}$; dado un \omterm{def:b3:topology:compact}{compacto}
numerable $K_j \subseteq [0, 1]$ con $K_j^{(j)} =
\{0\}$ y $K_j^{(j+1)} = \varnothing$, póngase
\[
K_{j+1} = \{0\} \cup \bigcup_{m \geq 1}
\Bigl(\frac1m + \varepsilon_m K_j\Bigr),
\]
con $\varepsilon_m$ suficientemente pequeño para que la $m$-ésima copia
esté en $\bigl(\frac1m, \frac1{m-1}\bigr)$. La derivación actúa copia a
copia (las copias viven en intervalos \omterm{def:b3:topology:topology}{abiertos} disjuntos), de modo que
$K_{j+1}^{(i)} = \{0\} \cup \bigcup_m\bigl(\frac1m +
\varepsilon_mK_j^{(i)}\bigr)$ para $i \leq j$; en $i = j$ esto se lee
$\{0\} \cup \{\frac1m\}$, y después $K_{j+1}^{(j+1)} =
\{0\}$ y $K_{j+1}^{(j+2)} = \varnothing$. Todo rango finito se alcanza.
Un conjunto \omterm{prop:b3:galois:perfect}{perfecto} es el otro
extremo: $P' = P$, la derivación nunca se mueve --- y Cantor--Bendixson
(pregunta 22) dice precisamente que todo cerrado se descompone en un
núcleo \omterm{prop:b3:galois:perfect}{perfecto}, invisible a la
derivación, y un resto numerable que la derivación va devorando.
\end{solution}
